% !TEX program = pdflatex
\documentclass[11pt,a4paper]{article}

% =========================================================================
% CRONOGRAMA DEL PROYECTO
% Proyecto: Horus Sentinel
% Gestión de Proyectos TICS (EST315) -- Semestre 2026-II -- Team Sobrevivientes
% =========================================================================

% --- PAQUETES DE CONFIGURACIÓN ---
\usepackage[utf8]{inputenc}
\usepackage[T1]{fontenc}
\usepackage[spanish,es-tabla]{babel}
% Compatibilidad con babel-spanish reciente (flechas de TikZ y signo de porcentaje)
\AtBeginDocument{\shorthandoff{<>}\spanishplainpercent}
\usepackage{geometry}
\geometry{
    top=3.6cm,
    bottom=2.3cm,
    left=2.2cm,
    right=2.2cm,
    % PARTE MODIFICABLE: headheight debe ser mayor que la altura del logo más alto (2.56cm = 73pt).
    headheight=80pt,
    headsep=15pt,
    footskip=20pt
}

\usepackage{graphicx}
\usepackage{fancyhdr}
\usepackage[table]{xcolor}
\usepackage{titlesec}
\usepackage{booktabs}
\usepackage{tabularx}
\usepackage{array}
\usepackage{amsmath}
\usepackage{enumitem}
\usepackage{microtype}
\usepackage{amssymb}
% --- RECURSOS VISUALES: GRÁFICOS Y CÓDIGO DE COLORES ---
\usepackage{tikz}
\usetikzlibrary{positioning,calc}
\definecolor{GrisFila}{RGB}{246, 248, 251}
\definecolor{CuadCerca}{RGB}{150, 30, 30}
\definecolor{CuadSatisf}{RGB}{180, 130, 20}
\definecolor{CuadInfo}{RGB}{10, 45, 90}
\definecolor{CuadMonit}{RGB}{110, 115, 125}
\newcommand{\etiqueta}[2]{{\setlength{\fboxsep}{2pt}\colorbox{#1}{\textcolor{white}{\scriptsize\bfseries #2}}}}
\newcommand{\cerca}{\etiqueta{CuadCerca}{Gestionar de cerca}}
\newcommand{\satisf}{\etiqueta{CuadSatisf}{Mantener satisfecho}}
\newcommand{\informado}{\etiqueta{CuadInfo}{Mantener informado}}
\newcommand{\monit}{\etiqueta{CuadMonit}{Monitorear}}
\usepackage{needspace}
\usepackage{longtable}
\definecolor{Rama1}{RGB}{10, 45, 90}
\definecolor{Rama2}{RGB}{40, 95, 150}
\definecolor{Rama3}{RGB}{180, 130, 20}
\definecolor{Rama4}{RGB}{30, 110, 60}
\definecolor{Rama5}{RGB}{110, 60, 140}
\definecolor{RaciR}{RGB}{150, 30, 30}
\definecolor{RaciA}{RGB}{10, 45, 90}
\definecolor{RaciC}{RGB}{180, 130, 20}
\definecolor{RaciI}{RGB}{140, 145, 155}
\definecolor{NivAlto}{RGB}{30, 110, 60}
\definecolor{NivMedio}{RGB}{180, 130, 20}
\definecolor{NivBajo}{RGB}{150, 30, 30}
\usetikzlibrary{patterns}

\usepackage{hyperref}

% Logos: funciona a 2 o 3 niveles de profundidad dentro de docs/latex/
\graphicspath{{../../assets/img/}{../../../assets/img/}}

% --- PALETA INSTITUCIONAL FIEI / UNA PUNO ---
\definecolor{AzulUNA}{RGB}{10, 45, 90}
\definecolor{DoradoFIEI}{RGB}{180, 130, 20}
\definecolor{GrisOscuro}{RGB}{50, 50, 50}
\definecolor{AzulClaro}{RGB}{228, 236, 246}

% --- CONFIGURACIÓN DE LISTAS ---
\setlist{topsep=2pt, itemsep=2pt, parsep=1pt}

% --- COLUMNAS DE TABLA ---
\newcolumntype{L}{>{\raggedright\arraybackslash}X}
\newcolumntype{P}[1]{>{\raggedright\arraybackslash}p{#1}}
\newcolumntype{C}[1]{>{\centering\arraybackslash}p{#1}}
\newcolumntype{Q}[1]{>{\raggedleft\arraybackslash}p{#1}}
\newcolumntype{R}{>{\raggedleft\arraybackslash}X}
\renewcommand{\arraystretch}{1.25}

% --- FORMATO DE ENLACES ---
\hypersetup{
    colorlinks=true,
    linkcolor=AzulUNA,
    urlcolor=AzulUNA,
    citecolor=AzulUNA,
    pdftitle={CRONOGRAMA DEL PROYECTO -- Horus Sentinel},
    pdfauthor={Team Sobrevivientes}
}

% --- FORMATO DE TÍTULOS ---
\titleformat{\section}{\color{AzulUNA}\normalfont\large\bfseries}{\thesection.}{0.7em}{}[\color{AzulUNA}\titlerule]
\titleformat{\subsection}{\color{AzulUNA}\normalfont\normalsize\bfseries}{\thesubsection.}{0.5em}{}
\titlespacing*{\section}{0pt}{9pt}{4pt}
\titlespacing*{\subsection}{0pt}{6pt}{2pt}

% --- ENCABEZADO Y PIE DE PÁGINA (fancyhdr) ---
\pagestyle{fancy}
\fancyhf{}

% ===== TAMAÑO DE LOS LOGOS (PARTE MODIFICABLE) =====
% logo_fiei.png  : height=1.45cm (lado izquierdo)
% logo_equipo.png: height=2.56cm (lado derecho; 25 % más grande porque la imagen tiene más borde blanco)
% Para cambiar el tamaño, modifique solo el valor height=... de cada \includegraphics.
% Si el logo más alto supera 2.6cm, aumente también headheight y top en la configuración de página.
\fancyhead[L]{%
    \includegraphics[height=1.45cm,keepaspectratio]{logo_fiei.png}%
}
\fancyhead[C]{%
    \scriptsize \textbf{\textcolor{AzulUNA}{UNIVERSIDAD NACIONAL DEL ALTIPLANO DE PUNO}}\\[2pt]
    \tiny \textcolor{GrisOscuro}{FACULTAD DE INGENIERÍA ESTADÍSTICA E INFORMÁTICA}\\[2pt]
    \scriptsize \textbf{\textcolor{AzulUNA}{Gestión de Proyectos TICS (EST315) \textbar\ Team Sobrevivientes}}\\[2pt]
    \footnotesize \textbf{\textcolor{DoradoFIEI}{CRONOGRAMA DEL PROYECTO}}%
}
\fancyhead[R]{%
    \includegraphics[height=2.56cm,keepaspectratio]{logo_equipo.png}%
}

\renewcommand{\headrulewidth}{0pt}
\renewcommand{\headrule}{}

\fancyfoot[L]{\scriptsize \textcolor{GrisOscuro}{Horus Sentinel}}
\fancyfoot[C]{\small \thepage}
\fancyfoot[R]{\scriptsize \textcolor{GrisOscuro}{Puno, Perú -- Semestre 2026-II}}
\renewcommand{\footrulewidth}{0.4pt}

\begin{document}

\begin{center}
    {\Large \textbf{\textcolor{AzulUNA}{CRONOGRAMA DEL PROYECTO}}}\\[3pt]
    {\normalsize \textbf{Horus Sentinel}}\\[4pt]
    \footnotesize \textbf{Fecha:} viernes 18/09/2026 (semana 4) \quad \textbar\quad \textbf{Versión:} 1.0 \quad \textbar\quad \textbf{Fase:} Planificación
\end{center}

\vspace{0.1cm}

\section{Datos Generales}
\noindent\begin{tabularx}{\textwidth}{@{}P{4.2cm} L@{}}
\toprule
\textbf{Proyecto:} & Horus Sentinel \\
\textbf{Responsable:} & Ruth Karina Apaza Solis -- Analista de Planificación \\
\textbf{Proceso PM4DEV:} & Gestión del cronograma \\
\textbf{Insumos:} & EDT y diccionario de la EDT, plan del equipo y Acta de Constitución \\
\textbf{Dedicación:} & 8 horas semanales por integrante (32 horas del equipo por semana) \\
\bottomrule
\end{tabularx}

\section{Método}
\begin{itemize}[leftmargin=*]
    \item \textbf{Actividades:} cada paquete de trabajo de la EDT es una actividad del cronograma. Las duraciones están en semanas de calendario, con la dedicación parcial del equipo.
    \item \textbf{Estimación PERT:} para cada actividad se estima una duración optimista ($o = 0.8\,m$), más probable ($m$) y pesimista ($p = 1.5\,m$). La duración esperada es $t_e = (o + 4m + p)/6$ y la desviación es $\sigma = (p - o)/6$. Las actividades terminadas usan su duración real.
    \item \textbf{Ruta crítica:} se calcula con inicio y fin tempranos (pasada hacia adelante) y tardíos (pasada hacia atrás). Las actividades con holgura cero forman la ruta crítica.
    \item \textbf{Calendario:} la semana 1 del semestre no se cuenta. En las semanas 9 (exámenes parciales) y 17 (exámenes finales) solo avanzan las actividades de gestión y documentación; el desarrollo, las pruebas y el trabajo de campo se detienen.
\end{itemize}

\section{Calendario del Proyecto}
{\small
\setlength{\tabcolsep}{4pt}
\noindent\begin{tabularx}{\textwidth}{@{}C{1.0cm} L C{1.0cm} L C{1.0cm} L@{}}
\toprule
\rowcolor{AzulClaro}
\textbf{Sem.} & \textbf{Fechas} & \textbf{Sem.} & \textbf{Fechas} & \textbf{Sem.} & \textbf{Fechas} \\
\midrule
2 & 31/08 -- 04/09 & 9 & 19/10 -- 23/10 \textcolor{RaciR}{\scriptsize(Exámenes parciales)} & 16 & 07/12 -- 11/12 \\
3 & 07/09 -- 11/09 & 10 & 26/10 -- 30/10 & 17 & 14/12 -- 18/12 \textcolor{RaciR}{\scriptsize(Exámenes finales)} \\
4 & 14/09 -- 18/09 & 11 & 02/11 -- 06/11 & 18 & \textcolor{GrisOscuro}{21/12 -- 25/12 \textcolor{RaciR}{\scriptsize(Cierre del semestre (24/12))}} \\
5 & 21/09 -- 25/09 & 12 & 09/11 -- 13/11 & 19 & \textcolor{GrisOscuro}{28/12 -- 01/01} \\
6 & 28/09 -- 02/10 & 13 & 16/11 -- 20/11 & 20 & \textcolor{GrisOscuro}{04/01 -- 08/01} \\
7 & 05/10 -- 09/10 & 14 & 23/11 -- 27/11 &  &  \\
8 & 12/10 -- 16/10 & 15 & 30/11 -- 04/12 &  &  \\
\bottomrule
\end{tabularx}
}

\Needspace{14\baselineskip}
\section{Red de Actividades}
Cada nodo muestra el código de la actividad y su duración esperada en semanas. Los nodos se ordenan de izquierda a derecha según sus predecesoras y los de borde rojo forman la ruta crítica.

\begin{center}
\begin{tikzpicture}[font=\tiny, nd/.style={draw=#1, fill=#1!10, line width=0.6pt, rounded corners=1.5pt, minimum width=1.0cm, minimum height=0.62cm, align=center, inner sep=1pt}, >=stealth]
\node[nd=Rama1] (n1_1) at (1.75,-1.05) {\textbf{1.1}\\ $t_e$ 1.00};
\node[nd=Rama1] (n1_2) at (2.91,-1.05) {\textbf{1.2}\\ $t_e$ 1.05};
\node[nd=Rama1] (n1_3) at (4.06,-1.05) {\textbf{1.3}\\ $t_e$ 1.05};
\node[nd=Rama1] (n1_5) at (8.68,-1.05) {\textbf{1.5}\\ $t_e$ 3.15};
\node[nd=RaciR, line width=1.2pt] (n2_1) at (0.60,0.00) {\textbf{2.1}\\ $t_e$ 1.00};
\node[nd=RaciR, line width=1.2pt] (n2_2) at (2.91,0.00) {\textbf{2.2}\\ $t_e$ 1.05};
\node[nd=RaciR, line width=1.2pt] (n2_3) at (1.75,0.00) {\textbf{2.3}\\ $t_e$ 1.57};
\node[nd=RaciR, line width=1.2pt] (n3_1) at (4.06,0.00) {\textbf{3.1}\\ $t_e$ 1.05};
\node[nd=Rama3] (n3_2) at (4.06,-2.10) {\textbf{3.2}\\ $t_e$ 1.57};
\node[nd=RaciR, line width=1.2pt] (n3_3) at (5.22,0.00) {\textbf{3.3}\\ $t_e$ 1.57};
\node[nd=RaciR, line width=1.2pt] (n3_4) at (6.37,0.00) {\textbf{3.4}\\ $t_e$ 2.10};
\node[nd=Rama3] (n3_5) at (6.37,-1.05) {\textbf{3.5}\\ $t_e$ 1.05};
\node[nd=RaciR, line width=1.2pt] (n3_6) at (7.52,0.00) {\textbf{3.6}\\ $t_e$ 2.10};
\node[nd=RaciR, line width=1.2pt] (n3_7) at (9.83,0.00) {\textbf{3.7}\\ $t_e$ 1.05};
\node[nd=Rama4] (n4_1) at (7.52,-1.05) {\textbf{4.1}\\ $t_e$ 1.05};
\node[nd=RaciR, line width=1.2pt] (n4_2) at (8.68,0.00) {\textbf{4.2}\\ $t_e$ 1.57};
\node[nd=RaciR, line width=1.2pt] (n4_3) at (10.98,0.00) {\textbf{4.3}\\ $t_e$ 1.05};
\node[nd=RaciR, line width=1.2pt] (n4_4) at (12.14,0.00) {\textbf{4.4}\\ $t_e$ 1.05};
\node[nd=RaciR, line width=1.2pt] (n5_1) at (13.29,0.00) {\textbf{5.1}\\ $t_e$ 1.05};
\node[nd=RaciR, line width=1.2pt] (n5_2) at (13.29,-1.05) {\textbf{5.2}\\ $t_e$ 1.05};
\node[nd=RaciR, line width=1.2pt] (n5_3) at (14.45,0.00) {\textbf{5.3}\\ $t_e$ 0.53};
\node[draw=AzulUNA, fill=AzulUNA, text=white, circle, inner sep=1pt, font=\tiny\bfseries] (ini) at (-0.55,0) {I};
\node[draw=AzulUNA, fill=AzulUNA, text=white, circle, inner sep=1pt, font=\tiny\bfseries] (fin) at (15.40,0) {F};
\draw[->, gray!60] (n2_1.east) to[out=0,in=180] (n1_1.west);
\draw[->, gray!60] (n1_1.east) to[out=0,in=180] (n1_2.west);
\draw[->, gray!60] (n1_2.east) to[out=0,in=180] (n1_3.west);
\draw[->, gray!60] (n4_1.east) to[out=0,in=180] (n1_5.west);
\draw[->, RaciR, thick] (ini) to[out=0,in=180] (n2_1.west);
\draw[->, RaciR, thick] (n2_3.east) to[out=0,in=180] (n2_2.west);
\draw[->, gray!60] (n1_1.east) to[out=0,in=180] (n2_2.west);
\draw[->, RaciR, thick] (n2_1.east) to[out=0,in=180] (n2_3.west);
\draw[->, RaciR, thick] (n2_2.east) to[out=0,in=180] (n3_1.west);
\draw[->, gray!60] (n2_2.east) to[out=0,in=180] (n3_2.west);
\draw[->, RaciR, thick] (n3_1.east) to[out=0,in=180] (n3_3.west);
\draw[->, gray!60] (n3_2.east) to[out=0,in=180] (n3_4.west);
\draw[->, RaciR, thick] (n3_3.east) to[out=0,in=180] (n3_4.west);
\draw[->, gray!60] (n3_3.east) to[out=0,in=180] (n3_5.west);
\draw[->, RaciR, thick] (n3_4.east) to[out=0,in=180] (n3_6.west);
\draw[->, RaciR, thick] (n4_2.east) to[out=0,in=180] (n3_7.west);
\draw[->, gray!60] (n3_4.east) to[out=0,in=180] (n4_1.west);
\draw[->, RaciR, thick] (n3_6.east) to[out=0,in=180] (n4_2.west);
\draw[->, RaciR, thick] (n3_7.east) to[out=0,in=180] (n4_3.west);
\draw[->, gray!60] (n3_5.east) to[out=0,in=180] (n4_3.west);
\draw[->, gray!60] (n4_1.east) to[out=0,in=180] (n4_3.west);
\draw[->, RaciR, thick] (n4_3.east) to[out=0,in=180] (n4_4.west);
\draw[->, RaciR, thick] (n4_4.east) to[out=0,in=180] (n5_1.west);
\draw[->, RaciR, thick] (n4_4.east) to[out=0,in=180] (n5_2.west);
\draw[->, RaciR, thick] (n5_1.east) to[out=0,in=180] (n5_3.west);
\draw[->, RaciR, thick] (n5_2.east) to[out=0,in=180] (n5_3.west);
\draw[->, gray!60] (n1_5.east) to[out=0,in=180] (n5_3.west);
\draw[->, gray!70] (n1_3.east) to[out=0,in=180] (fin);
\draw[->, RaciR, thick] (n5_3.east) to[out=0,in=180] (fin);
\end{tikzpicture}
\end{center}

\section{Estimación PERT}
Duraciones en semanas. Las filas resaltadas pertenecen a la ruta crítica.

\vspace{0.1cm}
{\small
\setlength{\tabcolsep}{3pt}
\begin{longtable}{@{}P{0.75cm} >{\raggedright\arraybackslash}p{4.3cm} P{1.95cm} Q{0.85cm} Q{0.85cm} Q{0.85cm} Q{0.95cm} Q{0.85cm} C{1.75cm}@{}}
\toprule
\rowcolor{AzulClaro}
\textbf{Cód.} & \textbf{Actividad} & \textbf{Pred.} & \textbf{$o$} & \textbf{$m$} & \textbf{$p$} & \textbf{$t_e$} & \textbf{$\sigma$} & \textbf{Estado} \\
\midrule
\endfirsthead
\toprule
\rowcolor{AzulClaro}
\textbf{Cód.} & \textbf{Actividad} & \textbf{Pred.} & \textbf{$o$} & \textbf{$m$} & \textbf{$p$} & \textbf{$t_e$} & \textbf{$\sigma$} & \textbf{Estado} \\
\midrule
\endhead
\bottomrule
\endfoot
1.1 & Viabilidad y Acta de Constitución & 2.1 & 1.00 & 1.00 & 1.00 & \textbf{1.00} & 0.00 & Terminada \\
1.2 & Planificación base & 1.1 & 0.80 & 1.00 & 1.50 & \textbf{1.05} & 0.12 & En curso \\
1.3 & Planes de soporte & 1.2 & 0.80 & 1.00 & 1.50 & \textbf{1.05} & 0.12 & Por iniciar \\
1.4 & Seguimiento y control & 1.3 & \multicolumn{5}{c}{\footnotesize nivel de esfuerzo: dura todo el proyecto} & En curso \\
1.5 & Convenio municipal & 4.1 & 2.40 & 3.00 & 4.50 & \textbf{3.15} & 0.35 & Por iniciar \\
\rowcolor{RaciR!8}\textbf{2.1} & Ficha técnica y requerimientos & Inicio & 1.00 & 1.00 & 1.00 & \textbf{1.00} & 0.00 & Terminada \\
\rowcolor{RaciR!8}\textbf{2.2} & Modelo de datos geoespacial & 2.3, 1.1 & 0.80 & 1.00 & 1.50 & \textbf{1.05} & 0.12 & En curso \\
\rowcolor{RaciR!8}\textbf{2.3} & Capacitación geoespacial & 2.1 & 1.20 & 1.50 & 2.25 & \textbf{1.57} & 0.17 & En curso \\
\rowcolor{RaciR!8}\textbf{3.1} & Datos simulados del cercado & 2.2 & 0.80 & 1.00 & 1.50 & \textbf{1.05} & 0.12 & Por iniciar \\
3.2 & Grafo de calles & 2.2 & 1.20 & 1.50 & 2.25 & \textbf{1.57} & 0.17 & Por iniciar \\
\rowcolor{RaciR!8}\textbf{3.3} & Mapa de calor (KDE) & 3.1 & 1.20 & 1.50 & 2.25 & \textbf{1.57} & 0.17 & Por iniciar \\
\rowcolor{RaciR!8}\textbf{3.4} & Ruta segura & 3.2, 3.3 & 1.60 & 2.00 & 3.00 & \textbf{2.10} & 0.23 & Por iniciar \\
3.5 & Capas por horario & 3.3 & 0.80 & 1.00 & 1.50 & \textbf{1.05} & 0.12 & Por iniciar \\
\rowcolor{RaciR!8}\textbf{3.6} & Reporte ciudadano & 3.4 & 1.60 & 2.00 & 3.00 & \textbf{2.10} & 0.23 & Por iniciar \\
\rowcolor{RaciR!8}\textbf{3.7} & Zonas anómalas & 4.2 & 0.80 & 1.00 & 1.50 & \textbf{1.05} & 0.12 & Por iniciar \\
4.1 & Pruebas del prototipo & 3.4 & 0.80 & 1.00 & 1.50 & \textbf{1.05} & 0.12 & Por iniciar \\
\rowcolor{RaciR!8}\textbf{4.2} & Captación de reportes & 3.6 & 1.20 & 1.50 & 2.25 & \textbf{1.57} & 0.17 & Por iniciar \\
\rowcolor{RaciR!8}\textbf{4.3} & Piloto de la ruta segura & 3.7, 3.5, 4.1 & 0.80 & 1.00 & 1.50 & \textbf{1.05} & 0.12 & Por iniciar \\
\rowcolor{RaciR!8}\textbf{4.4} & Pruebas de aceptación & 4.3 & 0.80 & 1.00 & 1.50 & \textbf{1.05} & 0.12 & Por iniciar \\
\rowcolor{RaciR!8}\textbf{5.1} & Manuales & 4.4 & 0.80 & 1.00 & 1.50 & \textbf{1.05} & 0.12 & Por iniciar \\
\rowcolor{RaciR!8}\textbf{5.2} & Informe final & 4.4 & 0.80 & 1.00 & 1.50 & \textbf{1.05} & 0.12 & Por iniciar \\
\rowcolor{RaciR!8}\textbf{5.3} & Acta de cierre & 5.1, 5.2, 1.5 & 0.40 & 0.50 & 0.75 & \textbf{0.53} & 0.06 & Por iniciar \\
\end{longtable}
}

\Needspace{12\baselineskip}
\section{Ruta Crítica}
{\small
\setlength{\tabcolsep}{3pt}
\begin{longtable}{@{}P{0.75cm} >{\raggedright\arraybackslash}p{5.2cm} C{1.6cm} C{1.6cm} C{1.6cm} C{1.6cm} Q{1.3cm}@{}}
\toprule
\rowcolor{AzulClaro}
\textbf{Cód.} & \textbf{Actividad} & \textbf{Inicio temprano} & \textbf{Fin temprano} & \textbf{Inicio tardío} & \textbf{Fin tardío} & \textbf{Holgura (sem.)} \\
\midrule
\endfirsthead
\toprule
\rowcolor{AzulClaro}
\textbf{Cód.} & \textbf{Actividad} & \textbf{Inicio temprano} & \textbf{Fin temprano} & \textbf{Inicio tardío} & \textbf{Fin tardío} & \textbf{Holgura (sem.)} \\
\midrule
\endhead
\bottomrule
\endfoot
1.1 & Viabilidad y Acta de Constitución & 07/09 & 11/09 & 09/09 & 16/09 & 0.58 \\
1.2 & Planificación base & 14/09 & 21/09 & 24/12 & 31/12 & 14.65 \\
1.3 & Planes de soporte & 21/09 & 28/09 & 31/12 & 07/01 & 14.65 \\
1.4 & Seguimiento y control & 28/09 & 07/01 & 28/09 & 07/01 & -- \\
1.5 & Convenio municipal & 11/11 & 02/12 & 14/12 & 05/01 & 4.68 \\
\rowcolor{RaciR!8}\textbf{2.1} & Ficha técnica y requerimientos & 31/08 & 04/09 & 31/08 & 04/09 & \textcolor{RaciR}{\textbf{0.00}} \\
\rowcolor{RaciR!8}\textbf{2.2} & Modelo de datos geoespacial & 16/09 & 24/09 & 16/09 & 24/09 & \textcolor{RaciR}{\textbf{0.00}} \\
\rowcolor{RaciR!8}\textbf{2.3} & Capacitación geoespacial & 07/09 & 16/09 & 07/09 & 16/09 & \textcolor{RaciR}{\textbf{0.00}} \\
\rowcolor{RaciR!8}\textbf{3.1} & Datos simulados del cercado & 24/09 & 01/10 & 24/09 & 01/10 & \textcolor{RaciR}{\textbf{0.00}} \\
3.2 & Grafo de calles & 24/09 & 05/10 & 01/10 & 13/10 & 1.05 \\
\rowcolor{RaciR!8}\textbf{3.3} & Mapa de calor (KDE) & 01/10 & 13/10 & 01/10 & 13/10 & \textcolor{RaciR}{\textbf{0.00}} \\
\rowcolor{RaciR!8}\textbf{3.4} & Ruta segura & 13/10 & 03/11 & 13/10 & 03/11 & \textcolor{RaciR}{\textbf{0.00}} \\
3.5 & Capas por horario & 13/10 & 27/10 & 30/11 & 07/12 & 6.77 \\
\rowcolor{RaciR!8}\textbf{3.6} & Reporte ciudadano & 03/11 & 18/11 & 03/11 & 18/11 & \textcolor{RaciR}{\textbf{0.00}} \\
\rowcolor{RaciR!8}\textbf{3.7} & Zonas anómalas & 30/11 & 07/12 & 30/11 & 07/12 & \textcolor{RaciR}{\textbf{0.00}} \\
4.1 & Pruebas del prototipo & 03/11 & 10/11 & 30/11 & 07/12 & 3.67 \\
\rowcolor{RaciR!8}\textbf{4.2} & Captación de reportes & 18/11 & 30/11 & 18/11 & 30/11 & \textcolor{RaciR}{\textbf{0.00}} \\
\rowcolor{RaciR!8}\textbf{4.3} & Piloto de la ruta segura & 07/12 & 21/12 & 07/12 & 21/12 & \textcolor{RaciR}{\textbf{0.00}} \\
\rowcolor{RaciR!8}\textbf{4.4} & Pruebas de aceptación & 21/12 & 28/12 & 21/12 & 28/12 & \textcolor{RaciR}{\textbf{0.00}} \\
\rowcolor{RaciR!8}\textbf{5.1} & Manuales & 28/12 & 05/01 & 28/12 & 05/01 & \textcolor{RaciR}{\textbf{0.00}} \\
\rowcolor{RaciR!8}\textbf{5.2} & Informe final & 28/12 & 05/01 & 28/12 & 05/01 & \textcolor{RaciR}{\textbf{0.00}} \\
\rowcolor{RaciR!8}\textbf{5.3} & Acta de cierre & 05/01 & 07/01 & 05/01 & 07/01 & \textcolor{RaciR}{\textbf{0.00}} \\
\end{longtable}
}

\vspace{0.1cm}
\noindent\textbf{Actividades críticas (holgura cero):} 2.1, 2.3, 2.2, 3.1, 3.3, 3.4, 3.6, 4.2, 3.7, 4.3, 4.4, 5.1, 5.2, 5.3.

\vspace{0.1cm}
\noindent La ruta crítica dura \textbf{18.75 semanas} de calendario: el proyecto requiere las semanas 2 a 20, es decir \textbf{19 semanas}, y termina el \textbf{07/01/2027}. El semestre ofrece 16 semanas (2 a 17), por lo que el cronograma excede el plazo en 3 semanas.

\begin{itemize}[leftmargin=*]
    \item Desviación estándar de la ruta crítica: $\sigma = 0.55$ semanas.
    \item Probabilidad de terminar hasta el 18/12/2026 (fin de la semana 17): $Z = (18.00 - 20.75)/0.55 = -5.03$, es decir \textbf{menos de 0.1\,\%}.
    \item Causas: la capacitación geoespacial antecede al modelo de datos; el mapa de calor, la ruta segura y el reporte ciudadano son secuenciales; la captación de reportes necesita la app terminada; y las semanas de exámenes detienen el desarrollo y el trabajo de campo.
    \item El resultado se eleva al equipo, como sponsor, para decidir la continuidad del proyecto.
\end{itemize}

\Needspace{20\baselineskip}
\section{Diagrama de Gantt por Semanas}
\begin{center}
\begin{tikzpicture}[font=\tiny]
\fill[AzulClaro] (0,0.05) rectangle (10.30,0.55);
\node[font=\tiny\bfseries, text=AzulUNA] at (0.27,0.30) {2};
\node[font=\tiny\bfseries, text=AzulUNA] at (0.81,0.30) {3};
\node[font=\tiny\bfseries, text=AzulUNA] at (1.36,0.30) {4};
\node[font=\tiny\bfseries, text=AzulUNA] at (1.90,0.30) {5};
\node[font=\tiny\bfseries, text=AzulUNA] at (2.44,0.30) {6};
\node[font=\tiny\bfseries, text=AzulUNA] at (2.98,0.30) {7};
\node[font=\tiny\bfseries, text=AzulUNA] at (3.52,0.30) {8};
\node[font=\tiny\bfseries, text=AzulUNA] at (4.07,0.30) {9};
\node[font=\tiny\bfseries, text=AzulUNA] at (4.61,0.30) {10};
\node[font=\tiny\bfseries, text=AzulUNA] at (5.15,0.30) {11};
\node[font=\tiny\bfseries, text=AzulUNA] at (5.69,0.30) {12};
\node[font=\tiny\bfseries, text=AzulUNA] at (6.23,0.30) {13};
\node[font=\tiny\bfseries, text=AzulUNA] at (6.78,0.30) {14};
\node[font=\tiny\bfseries, text=AzulUNA] at (7.32,0.30) {15};
\node[font=\tiny\bfseries, text=AzulUNA] at (7.86,0.30) {16};
\node[font=\tiny\bfseries, text=AzulUNA] at (8.40,0.30) {17};
\node[font=\tiny\bfseries, text=AzulUNA] at (8.94,0.30) {18};
\node[font=\tiny\bfseries, text=AzulUNA] at (9.49,0.30) {19};
\node[font=\tiny\bfseries, text=AzulUNA] at (10.03,0.30) {20};
\fill[GrisFila!60!gray!25] (3.79,0.05) rectangle (4.34,-11.04);
\fill[GrisFila!60!gray!25] (8.13,0.05) rectangle (8.67,-11.04);
\fill[RaciR!6] (8.67,0.05) rectangle (10.30,-11.04);
\draw[gray!25] (0.00,0.05) -- (0.00,-11.04);
\draw[gray!25] (0.54,0.05) -- (0.54,-11.04);
\draw[gray!25] (1.08,0.05) -- (1.08,-11.04);
\draw[gray!25] (1.63,0.05) -- (1.63,-11.04);
\draw[gray!25] (2.17,0.05) -- (2.17,-11.04);
\draw[gray!25] (2.71,0.05) -- (2.71,-11.04);
\draw[gray!25] (3.25,0.05) -- (3.25,-11.04);
\draw[gray!25] (3.79,0.05) -- (3.79,-11.04);
\draw[gray!25] (4.34,0.05) -- (4.34,-11.04);
\draw[gray!25] (4.88,0.05) -- (4.88,-11.04);
\draw[gray!25] (5.42,0.05) -- (5.42,-11.04);
\draw[gray!25] (5.96,0.05) -- (5.96,-11.04);
\draw[gray!25] (6.51,0.05) -- (6.51,-11.04);
\draw[gray!25] (7.05,0.05) -- (7.05,-11.04);
\draw[gray!25] (7.59,0.05) -- (7.59,-11.04);
\draw[gray!25] (8.13,0.05) -- (8.13,-11.04);
\draw[gray!25] (8.67,0.05) -- (8.67,-11.04);
\draw[gray!25] (9.22,0.05) -- (9.22,-11.04);
\draw[gray!25] (9.76,0.05) -- (9.76,-11.04);
\draw[gray!25] (10.30,0.05) -- (10.30,-11.04);
\node[anchor=east, font=\tiny\bfseries, text=Rama1] at (-0.1,-0.40) {1 Gestión del proyecto};
\node[anchor=east, font=\tiny, text width=4.7cm, align=right] at (-0.1,-0.80) {1.1 Viabilidad y Acta de Constitución};
\fill[Rama1] (0.54,-0.92) rectangle (1.08,-0.68);
\draw[white, line width=0.4pt] (0.54,-0.80) -- (1.08,-0.80);
\node[anchor=east, font=\tiny, text width=4.7cm, align=right] at (-0.1,-1.20) {1.2 Planificación base};
\fill[Rama1] (1.08,-1.32) rectangle (1.65,-1.08);
\node[anchor=east, font=\tiny, text width=4.7cm, align=right] at (-0.1,-1.60) {1.3 Planes de soporte};
\fill[Rama1] (1.65,-1.72) rectangle (2.22,-1.48);
\node[anchor=east, font=\tiny, text width=4.7cm, align=right] at (-0.1,-2.00) {1.4 Seguimiento y control};
\fill[pattern=north east lines, pattern color=GrisOscuro!60] (2.22,-2.11) rectangle (10.16,-1.89);
\draw[GrisOscuro!60] (2.22,-2.11) rectangle (10.16,-1.89);
\node[anchor=east, font=\tiny, text width=4.7cm, align=right] at (-0.1,-2.40) {1.5 Convenio municipal};
\fill[Rama1] (5.64,-2.52) rectangle (7.35,-2.28);
\node[anchor=east, font=\tiny\bfseries, text=Rama2] at (-0.1,-2.80) {2 Análisis y diseño};
\node[anchor=east, font=\tiny, text width=4.7cm, align=right] at (-0.1,-3.20) {2.1 Ficha técnica y requerimientos};
\fill[RaciR] (0.00,-3.32) rectangle (0.54,-3.08);
\draw[white, line width=0.4pt] (0.00,-3.20) -- (0.54,-3.20);
\node[anchor=east, font=\tiny, text width=4.7cm, align=right] at (-0.1,-3.60) {2.2 Modelo de datos geoespacial};
\fill[RaciR] (1.40,-3.72) rectangle (1.97,-3.48);
\node[anchor=east, font=\tiny, text width=4.7cm, align=right] at (-0.1,-4.00) {2.3 Capacitación geoespacial};
\fill[RaciR] (0.54,-4.12) rectangle (1.40,-3.88);
\node[anchor=east, font=\tiny\bfseries, text=Rama3] at (-0.1,-4.40) {3 Construcción del software};
\node[anchor=east, font=\tiny, text width=4.7cm, align=right] at (-0.1,-4.80) {3.1 Datos simulados del cercado};
\fill[RaciR] (1.97,-4.92) rectangle (2.53,-4.68);
\node[anchor=east, font=\tiny, text width=4.7cm, align=right] at (-0.1,-5.20) {3.2 Grafo de calles};
\fill[Rama3] (1.97,-5.32) rectangle (2.82,-5.08);
\node[anchor=east, font=\tiny, text width=4.7cm, align=right] at (-0.1,-5.60) {3.3 Mapa de calor (KDE)};
\fill[RaciR] (2.53,-5.72) rectangle (3.39,-5.48);
\node[anchor=east, font=\tiny, text width=4.7cm, align=right] at (-0.1,-6.00) {3.4 Ruta segura};
\fill[RaciR] (3.39,-6.12) rectangle (5.07,-5.88);
\node[anchor=east, font=\tiny, text width=4.7cm, align=right] at (-0.1,-6.40) {3.5 Capas por horario};
\fill[Rama3] (3.39,-6.52) rectangle (4.50,-6.28);
\node[anchor=east, font=\tiny, text width=4.7cm, align=right] at (-0.1,-6.80) {3.6 Reporte ciudadano};
\fill[RaciR] (5.07,-6.92) rectangle (6.21,-6.68);
\node[anchor=east, font=\tiny, text width=4.7cm, align=right] at (-0.1,-7.20) {3.7 Zonas anómalas};
\fill[RaciR] (7.06,-7.32) rectangle (7.63,-7.08);
\node[anchor=east, font=\tiny\bfseries, text=Rama4] at (-0.1,-7.60) {4 Pruebas y validación con clientes};
\node[anchor=east, font=\tiny, text width=4.7cm, align=right] at (-0.1,-8.00) {4.1 Pruebas del prototipo};
\fill[Rama4] (5.07,-8.12) rectangle (5.64,-7.88);
\node[anchor=east, font=\tiny, text width=4.7cm, align=right] at (-0.1,-8.40) {4.2 Captación de reportes};
\fill[RaciR] (6.21,-8.52) rectangle (7.06,-8.28);
\node[anchor=east, font=\tiny, text width=4.7cm, align=right] at (-0.1,-8.80) {4.3 Piloto de la ruta segura};
\fill[RaciR] (7.63,-8.92) rectangle (8.74,-8.68);
\node[anchor=east, font=\tiny, text width=4.7cm, align=right] at (-0.1,-9.20) {4.4 Pruebas de aceptación};
\fill[RaciR] (8.74,-9.32) rectangle (9.31,-9.08);
\node[anchor=east, font=\tiny\bfseries, text=Rama5] at (-0.1,-9.60) {5 Cierre};
\node[anchor=east, font=\tiny, text width=4.7cm, align=right] at (-0.1,-10.00) {5.1 Manuales};
\fill[RaciR] (9.31,-10.12) rectangle (9.88,-9.88);
\node[anchor=east, font=\tiny, text width=4.7cm, align=right] at (-0.1,-10.40) {5.2 Informe final};
\fill[RaciR] (9.31,-10.52) rectangle (9.88,-10.28);
\node[anchor=east, font=\tiny, text width=4.7cm, align=right] at (-0.1,-10.80) {5.3 Acta de cierre};
\fill[RaciR] (9.88,-10.92) rectangle (10.16,-10.68);
\draw[RaciR, dashed, thick] (8.67,0.6) -- (8.67,-11.04);
\node[anchor=south, text=RaciR, font=\tiny\bfseries] at (8.67,0.55) {fin de la semana 17};
\end{tikzpicture}
\\[2pt]
{\scriptsize \textcolor{RaciR}{\rule{9pt}{5pt}}~Ruta crítica \quad \textcolor{Rama1}{\rule{9pt}{5pt}}~Rama 1 \quad \textcolor{Rama2}{\rule{9pt}{5pt}}~Rama 2 \quad \textcolor{Rama3}{\rule{9pt}{5pt}}~Rama 3 \quad \textcolor{Rama4}{\rule{9pt}{5pt}}~Rama 4 \quad \textcolor{Rama5}{\rule{9pt}{5pt}}~Rama 5 \quad \tikz{\fill[pattern=north east lines, pattern color=GrisOscuro!60] (0,0) rectangle (0.32,0.18);}~Nivel de esfuerzo \quad \tikz{\fill[GrisFila!60!gray!25] (0,0) rectangle (0.32,0.18);}~Semana de exámenes}
\end{center}

\Needspace{12\baselineskip}
\section{Línea Base del Cronograma}
{\small
\rowcolors{2}{white}{GrisFila}
\noindent\begin{tabularx}{\textwidth}{@{}L C{0.9cm} C{2.75cm} C{2.75cm} C{2.3cm}@{}}
\toprule
\rowcolor{AzulClaro}
\textbf{Hito} & \textbf{Act.} & \textbf{Fecha del Acta} & \textbf{Línea base} & \textbf{Diferencia} \\
\midrule
Acta de Constitución aprobada & 1.1 & 11/09 (sem.\,3) & 11/09 (sem.\,3) & \textcolor{Rama4}{\textbf{a tiempo}} \\
Prototipo con datos simulados & 4.1 & 16/10 (sem.\,8) & 10/11 (sem.\,12) & \textcolor{RaciR}{\textbf{25 días después}} \\
Piloto con ciudadanos & 4.3 & 04/12 (sem.\,15) & 21/12 (sem.\,18) & \textcolor{RaciR}{\textbf{17 días después}} \\
MVP y convenio municipal & 4.4 & 11/12 (sem.\,16) & 28/12 (sem.\,19) & \textcolor{RaciR}{\textbf{17 días después}} \\
Acta de cierre & 5.3 & 18/12 (sem.\,17) & 07/01 (sem.\,20) & \textcolor{RaciR}{\textbf{20 días después}} \\
\bottomrule
\end{tabularx}
}

\vspace{0.2cm}
\noindent La línea base del cronograma se aprueba con esta versión. Durante la ejecución:
\begin{itemize}[leftmargin=*]
    \item Cada martes se registra el avance real de cada actividad y se compara con esta línea base (variación del cronograma e índice SPI).
    \item Un retraso mayor a 3 días en una actividad de la ruta crítica activa una acción correctiva en la reunión de control.
    \item Un cambio de fecha de un hito requiere una solicitud de cambio aprobada por el equipo.
\end{itemize}

\section{Firmas}
\noindent En señal de conformidad con el presente cronograma, los integrantes del equipo firman a continuación:

\vspace{1.5cm}
\begin{center}
\begin{tabular}{c c c}
    \parbox{6.5cm}{\centering \rule{6.0cm}{0.6pt}\\[5pt] \textbf{Clyde Neil Paricahua Pari}\\[2pt] {\footnotesize Project Manager}}
    & \hspace{0.8cm} &
    \parbox{6.5cm}{\centering \rule{6.0cm}{0.6pt}\\[5pt] \textbf{Ruth Karina Apaza Solis}\\[2pt] {\footnotesize Analista de Planificación}}
    \\[2.0cm]
    \parbox{6.5cm}{\centering \rule{6.0cm}{0.6pt}\\[5pt] \textbf{Eliseo Tarqui Ajahuana}\\[2pt] {\footnotesize Especialista Técnico / Arquitectura}}
    & \hspace{0.8cm} &
    \parbox{6.5cm}{\centering \rule{6.0cm}{0.6pt}\\[5pt] \textbf{Viviana Mary Quisbert Quispe}\\[2pt] {\footnotesize Gestora de Comunicación, Riesgos y Documentación}}
\end{tabular}
\end{center}
\end{document}
